\section{General framework}
\label{sec:framework}
This section sets out a framework that maps an identified macro-financial shock
into the probability of default (PD) of a credit portfolio. The framework is not
specific to any single shock: it applies to any innovation to a designated
variable of interest in a macro-financial system, such as a monetary,
financial-stress, oil-price, or uncertainty innovation. Geopolitical risk is the
application we develop in the empirical part of the paper. Geopolitical risk is
exogenous to the financial system and is transmitted through financial markets,
the real economy, and safety and security \citep{buch2024global}, which makes it
a natural application for a method designed to trace macro-financial shocks into
bank credit risk.
\medskip
The framework has two steps. In the first step we estimate, independently, two
blocks. A Gaussian vector autoregression (VAR) characterizes the joint dynamics
of the variable of interest and a set of macro-financial variables collected in
an $n$-dimensional vector $Y_t$. A Merton--Vasicek satellite model maps
macroeconomic conditions into the portfolio default probability through a latent
systematic factor $Z_t$, which we reconstruct from observed default rates and
relate to a subset of the VAR variables, denoted $Y_t^{(s)}$, by a Bayesian
linear regression. In the second step we identify the innovation to the variable
of interest, compute its generalized impulse response on $Y_t$
(\citealp{koop1996impulse}), propagate this response through the satellite to the
systematic factor, and map it into the portfolio default probability through a
closed-form expression that accounts for the nonlinearity of the Merton mapping.
\medskip
Two design choices are central to the framework and to its empirical relevance.
First, the satellite is estimated by Bayesian methods rather than by ordinary
least squares, so that estimation uncertainty in the credit-risk bridge is
quantified and propagated, coherently with the Bayesian VAR, into the default
probability responses. Second, because the VAR and the satellite are estimated
independently and are linked only through the impulse-response propagation of
Section~\ref{sec:irf_transmission}, the macro-financial sample used for the VAR
and the default sample used to reconstruct $Z_t$ need not share the same length,
start date, or frequency span. The VAR can be estimated on a long macro-financial
history, while $Z_t$ is reconstructed on the shorter default record. This
modularity matters in practice: bank default series are typically short, a few
years to at most two decades, relative to the macro-financial series, so a single
jointly estimated system would be infeasible or weakly identified. The nonlinear
Merton mapping further implies that linear impulse responses of the factor cannot
be transformed directly into default-probability responses; the closed-form
generalized impulse responses derived below remain valid under this nonlinearity
and, in the linear case, coincide with the ordering-invariant generalized impulse
responses of \citet{pesaranShin1998}. Figure~\ref{fig:twostep_approach}
summarizes the framework, and the following subsections give the details.
% --- In document ---
\begin{figure}[!ht]
\centering
\resizebox{0.96\textwidth}{!}{%
\begin{tikzpicture}[
  font=\small,
  >=Latex,
  box/.style={rounded corners=10pt, draw=none, text=white, align=center,
              minimum width=6.2cm, minimum height=1.25cm, inner sep=7pt},
  boxsm/.style={rounded corners=10pt, draw=none, text=white, align=center,
                minimum width=6.0cm, minimum height=1.1cm, inner sep=7pt},
  stepbadge/.style={draw=black!40, fill=white, rounded corners=3pt,
                    inner xsep=6pt, inner ysep=3pt, font=\small\bfseries},
  frame/.style={draw=gray!45, rounded corners=18pt, line width=0.9pt, inner sep=11pt},
  arrow/.style={-Latex, line width=1.15pt, draw=gray!65},
  darrow/.style={<->, line width=1.15pt, draw=gray!65}
]
% ==========================================
% Palette "beige / brown" (IRF-like theme)
% ==========================================
\definecolor{tanA}{RGB}{170,128,79}   % medium brown
\definecolor{tanB}{RGB}{200,161,111}  % light brown
\definecolor{tanC}{RGB}{232,216,193}  % beige (for frames)
\definecolor{tanG}{RGB}{55,55,55}     % dark gray (soft black)
% ------------------------------------------
% Layout parameters
% ------------------------------------------
\def\xcol{8.3cm}
\def\yrow{2.5cm}
\def\ysep{3.2cm}
\def\stepTwoShift{0cm}
% =========================
% STEP 1
% =========================
\node (maco) at (0,0) [box, fill=tanA] {Macro-financial sample (long)\\
$Y_t = (Y_{t,i})_{i\in I}$ with $I\subset\{1,\dots,n\}$};
\node (sysf) at (\xcol,0) [box, fill=tanB] {Default sample (short) $(d_t)_{t}$\\[1mm]
Realized systemic factor\\
$Z_t=\dfrac{\Phi^{-1}\!\left(\widehat p \right)-\sqrt{1-\rho}\,\Phi^{-1}\!\left(d_t\right)}{\sqrt{\rho}}$};
\node (svar) at (0,-\yrow) [boxsm, fill=tanA] {Bayesian VAR (reduced form)\\
$Y_t = c + \sum_{j=1}^{p} A_j Y_{t-j} + u_t,\quad u_t\sim\mathcal N(0,\Sigma_u)$};
\node (link) at (\xcol,-\yrow) [boxsm, fill=tanB] {Bayesian satellite linking $Z_t$ to the real economy\\
$Z_t = \beta_0 + \beta^\top Y_t^{(s)} + \eta_t$};
\draw[arrow] (maco.south) -- (svar.north);
\draw[arrow] (sysf.south) -- (link.north);
\draw[darrow] (svar.east) -- (link.west);
% Frame Step 1
\node[frame, draw=tanC!70!black] (frame1) [fit=(maco)(sysf)(svar)(link)] {};
\node[stepbadge] at ($(frame1.west)+(-1.5cm,0)$) {Step 1};
% =========================
% STEP 2
% =========================
\coordinate (step2anchor) at ($(frame1.center)+(0,-\ysep-\yrow)$);
\node (irfY) at ($(step2anchor)+(\stepTwoShift,0)$) [box, fill=tanG] {Generalized impulse response of the shock of interest on $Y_t$\\
$\psi_Y(h,\delta,\Omega_{t-1})=\mathbb{E}(Y_{t+h}\mid u_t=\delta,\Omega_{t-1})-\mathbb{E}(Y_{t+h}\mid \Omega_{t-1})$};
\node (irfZ) [box, fill=tanG, below=10mm of irfY] {Generalized impulse response of the shock of interest on $Z$\\
$\psi_Z(h,\delta,\Omega_{t-1})=\beta^\top\!\big(\mathbb{E}(Y_{t+h}\mid u_t=\delta,\Omega_{t-1})-\mathbb{E}(Y_{t+h}\mid \Omega_{t-1})\big)$};
\node (irfPD) [box, fill=tanG, below=10mm of irfZ] {Generalized impulse response of the shock of interest on the PD\\
$\psi_{PD}(h,\delta,\Omega_{t-1})=\mathbb{E}(\pi_{t+h}\mid u_t=\delta,\Omega_{t-1})-\mathbb{E}(\pi_{t+h}\mid \Omega_{t-1})$};
\draw[arrow] (irfY.south) -- (irfZ.north);
\draw[arrow] (irfZ.south) -- (irfPD.north);
\node[frame, draw=tanC!70!black] (frame2) [fit=(irfY)(irfZ)(irfPD)] {};
\node[stepbadge] at ($(frame2.west)+(-1.5cm,0)$) {Step 2};
% =========================
% Connector Step 1 -> Step 2
% =========================
\draw[arrow] (frame1.south) -- (frame2.north);
\end{tikzpicture}%
}
\caption{\textbf{A two-step approach}\\
\textit{Notes: The first step combines a macro-financial sample with a default
sample to estimate a Bayesian VAR and a Bayesian satellite model linking the
realized systemic factor to the real economy. The two blocks are estimated
independently, so the two samples need not have the same length. The second step
uses these models to assess the impact of the shock of interest (the
geopolitical-risk shock in our application) on macroeconomic variables, the
systemic factor, and portfolio PDs.}}
\label{fig:twostep_approach}
\end{figure}
\subsection{The VAR--Merton framework}\label{sec:var-merton framework}
The first step combines two components, estimated independently. A Gaussian VAR
characterizes the joint dynamics of the variable of interest and a set of
macro-financial variables. A Bayesian Merton--Vasicek satellite maps
macroeconomic conditions into the portfolio default probability. The two
components are connected in Section~\ref{sec:irf_transmission}.
\subsubsection{Macroeconomic dynamics}
Let $G_t$ denote the variable of interest (the geopolitical risk index in our
application) and let
$\mathbf{X}_t = (X_{1,t},\dots,X_{n-1,t})^\top \in \mathbb{R}^{n-1}$ collect the
remaining macro-financial variables. We order the variable of interest first and
stack the variables as
\[
    Y_t = (G_t,\ \mathbf{X}_t^\top)^\top \in \mathbb{R}^n.
\]
The ordering fixes notation; its role in identification is discussed in
Section~\ref{sec:irf_transmission}. The VAR($p$) model is
\begin{equation}
    Y_t = c + \sum_{i=1}^{p} A_i Y_{t-i} + u_t,\qquad
    u_t \stackrel{\text{i.i.d.}}{\sim} \mathcal N(\mathbf{0}_n, \Sigma_u),
    \quad t \in \mathbb{Z},
    \label{eq:var}
\end{equation}
with $c \in \mathbb{R}^n$, $A_i \in \mathbb{R}^{n \times n}$, and
$\Sigma_u \succ 0$. Following \citet{caldara2022geopolitical}, we estimate the
model under a Normal--Inverse--Wishart prior on $(A_1,\dots,A_p,c,\Sigma_u)$,
which combines an inverse--Wishart prior on $\Sigma_u$ with a conjugate Gaussian
prior on the coefficient matrices. We impose no prior restriction on the
coefficients of the $G$ equation, so that all lagged macro-financial variables
are allowed to affect the variable of interest.
\subsubsection{Default risk dynamics via a Merton satellite model}
We map macroeconomic conditions into default risk through a one-factor structure
based on \citet{merton1974pricing} and \citet{vasicek2002distribution}. We treat
the credit portfolio as an asymptotically granular and homogeneous pool of
obligors exposed to a single Gaussian systematic factor $Z_t$. Under this
asymptotic single-risk-factor assumption, all obligors share a common
unconditional probability of default $p \in (0,1)$ and a common asset
correlation $\rho \in [0,1)$, so the portfolio is summarized by a single
conditional default probability. The point-in-time probability of default
conditional on $Z_t$ is
\begin{equation}
    \pi(Z_t) = \Phi\!\left(
        \frac{\Phi^{-1}(p) - \sqrt{\rho}\, Z_t}{\sqrt{1-\rho}} \right),
    \label{eq:pd_conditional_z}
\end{equation}
where $\Phi(\cdot)$ is the standard normal cumulative distribution function. The
factor is oriented so that higher $Z_t$ corresponds to better systematic credit
conditions and a lower conditional default probability.
We reconstruct the factor and calibrate the parameters from the observed
portfolio default rates $\{d_1,\dots,d_T\}$, where $d_t \in (0,1)$ is the
portfolio default rate at time $t$ over an estimation sample of $T$ periods. This
sample is specific to the default block and may be considerably shorter than the
macro-financial sample underlying the VAR, since the two blocks communicate only
through the variables $Y^{(s)}_t$ and not through a shared estimation window. The
through-the-cycle probability of default is estimated by the sample mean of the
observed default rates,
\begin{equation}
    \widehat{p} = \frac{1}{T}\sum_{t=1}^T d_t.
    \label{eq:longrun_pd}
\end{equation}
We interpret the observed default rate as the realized conditional default
probability, $d_t = \pi(Z_t)$, which holds exactly in the asymptotically
granular limit and approximately for finite portfolios. Inverting
\eqref{eq:pd_conditional_z} for a given asset correlation $\rho$ yields the
implied factor
\begin{equation}
    Z_t(\rho,\widehat p)
    = \frac{\Phi^{-1}(\widehat p) - \sqrt{1-\rho}\,\Phi^{-1}(d_t)}{\sqrt{\rho}}.
    \label{eq:determines_z_hist}
\end{equation}
We select the asset correlation $\widehat\rho \in [0,1)$ such that the
reconstructed factor has unit unconditional variance,
$\Var_t[Z_t(\widehat\rho,\widehat p)] = 1$, and set
$Z_t := Z_t(\widehat\rho,\widehat p)$. This normalization fixes the scale of the
latent factor and ensures comparability with the standard Vasicek
parameterization.
The satellite model relates the reconstructed factor $Z_t$ to macroeconomic
conditions. If the Merton structure holds, the systematic component of credit
risk co-moves with observable macro-financial variables. Let
$Y^{(s)}_t \subseteq \{Y_t,\dots,Y_{t-q}\}$ denote the subset of variables and
lags entering the satellite equation. We specify
\begin{equation}
    Z_t = \beta_0 + \beta^\top Y^{(s)}_t + \eta_t,
    \qquad
    \eta_t \sim \mathcal{N}(0,\sigma_\eta^2),
    \qquad
    \E[\eta_t \mid Y^{(s)}_t] = 0.
    \label{eq:Z}
\end{equation}
We estimate \eqref{eq:Z} by Bayesian conjugate linear regression rather than by
ordinary least squares. Conditional on a set of regressors, we adopt the
non-informative prior $p(\beta_0,\beta,\sigma_\eta^2)\propto \sigma_\eta^{-2}$,
which yields a Normal--Inverse--Gamma posterior for $(\beta_0,\beta,\sigma_\eta^2)$
centered at the least-squares estimates and delivering the full posterior
distribution of the satellite parameters. The selection of the variables and
lags in $Y^{(s)}_t$ is described in Section~\ref{sec:lag_section}.
The Bayesian treatment of the satellite is motivated by three considerations.
First, default samples are short, so the satellite coefficients are estimated
imprecisely; the posterior quantifies this estimation uncertainty and, through
the propagation of Section~\ref{sec:irf_transmission}, carries it into the
default-probability responses as credible bands rather than point responses.
Second, it makes the two blocks coherent: the VAR is already estimated under a
Normal--Inverse--Wishart posterior, and a Bayesian satellite yields a single
posterior distribution for the PD response whose variance is decomposable into a
VAR-block contribution and a satellite-block contribution. Third, the posterior
of $\sigma_\eta^2$ enters the conditional variance of $Z_{t+h}$ in
Proposition~\ref{prop:Var}, so estimation uncertainty in the bridge affects both
the location and the dispersion of the implied default probabilities.
The satellite specification is itself uncertain. On a short default sample,
several choices of variables and lags in $Y^{(s)}_t$ fit comparably, and
conditioning inference on a single specification understates uncertainty. We
therefore consider a set of candidate specifications $\{s\}$ and assign each a
weight $w_s \propto \exp(-\tfrac{1}{2}\,\Delta\mathrm{AIC}_s)$, normalized over
the specifications retained up to a target cumulative weight. We report
model-averaged objects, obtained by mixing the posterior draws across
specifications in proportion to $w_s$. The reported coefficients, credible
intervals, and PD responses then integrate over parameter uncertainty within
each specification and over specification uncertainty across specifications. We
apply this averaging to each information set used in the empirical analysis.
\subsection{Shock transmission to default risk}\label{sec:irf_transmission}
This section derives the response of portfolio default probabilities to the shock
of interest (the geopolitical-risk shock in our application). Default
probabilities are obtained from the systematic factor through the bounded,
nonlinear Merton mapping \eqref{eq:pd_conditional_z}, so the linear impulse
response of the factor cannot be transformed directly into the response of the
default probability. We proceed in three steps: we compute the generalized
impulse response of the macro-financial variables to the innovation of interest
(\citealp{koop1996impulse}), propagate it to the systematic factor through the
satellite equation, and map it into default probabilities through a closed-form
expression that accounts for the nonlinearity of the Merton transformation. The
propositions are stated conditional on the model parameters; in estimation each
object is evaluated over the joint posterior draws of the VAR and the Bayesian
satellite, and reported as posterior medians with credible bands.
\subsubsection{Generalized impulse responses of macroeconomic variables}
We compute generalized impulse response functions (GIRFs) from the VAR in
\eqref{eq:var} and treat the reduced-form innovation to the equation of the
variable of interest as the shock. Isolating a single innovation does not require
ordering the remaining variables of $Y_t$, but it imposes one identifying
assumption: the innovation of interest is contemporaneously exogenous to the rest
of the system. By Proposition~3.1 of \citet{pesaranShin1998}, this assumption is
equivalent to a recursive (Cholesky) scheme in which the variable of interest is
ordered first, with the ordering of the remaining variables immaterial. We adopt
it following \citet{caldara2022geopolitical}, who construct the GPR index
independently of macro-financial outcomes.
We let $\Omega_{t-1}$ denote the information set at time $t-1$, the set of all
possible histories of the process up to $t-1$, with $\omega_{t-1} \in
\Omega_{t-1}$ a particular realization. Let $g \in \{1,\dots,n\}$ be the index of
the variable of interest in $Y_t$ and let $e_g \in \R^n$ denote the $g$-th
canonical unit vector. In our application this variable is ordered first, so
$g=1$. Its reduced-form innovation is $u_{gt} = e_g' u_t$.
\begin{defn}[\textbf{Generalized impulse response function}]\label{prop:GIRF}
    \textit{For a horizon $h \ge 0$, a history $\omega_{t-1} \in \Omega_{t-1}$,
    and a reduced-form shock vector $\delta \in \R^{n}$ applied to $u_t$,
    the generalized impulse response function (GIRF) of the vector $Y_t$
    is defined as}
    \begin{equation}
         \psi_Y(h,\delta,\omega_{t-1})
        := \E\!\big[Y_{t+h}\mid u_t=\delta,\ \Omega_{t-1} = \omega_{t-1}\big]
        - \E\!\big[Y_{t+h}\mid \Omega_{t-1} =\omega_{t-1}\big].
    \label{eq:GIRF-def}
    \end{equation}
\end{defn}
Rather than shocking all components of $u_t$, we focus on a shock to the single
innovation of the variable of interest.
\begin{defn}[\textbf{GIRF to a shock of interest}]\label{prop:GIRF-GPR}
    \textit{Let $g$ denote the index of the variable of interest in $Y_t$ and let
    $u_{gt} = e_g' u_t$ be its reduced-form innovation. For a horizon
    $h \ge 0$, a history $\omega_{t-1} \in \Omega_{t-1}$, and a scalar shock
    size $\delta_g \in \R$ applied to $u_{gt}$, the GIRF of $Y_t$ to this shock
    is}
    \begin{equation}
        \psi_Y^g(h,\delta_g,\omega_{t-1})
        := \E\!\big[Y_{t+h}\mid u_{gt}=\delta_g,\ \omega_{t-1}\big]
        - \E\!\big[Y_{t+h}\mid \omega_{t-1}\big].
    \label{eq:GIRF-def-GPR}
    \end{equation}
\end{defn}
Let $\{\Psi_h\}_{h \ge 0}$ denote the moving-average coefficients of the VAR, so
that
\begin{equation}
    Y_t = \mu + \sum_{h=0}^\infty \Psi_h u_{t-h}, \qquad \Psi_0 = I_n,
\end{equation}
and let $\Sigma_u = \Var(u_t)$ with $(g,g)$-element $\sigma_{gg}$. In a linear
Gaussian VAR the GIRF does not depend on the history $\omega_{t-1}$, and
\citet{pesaranShin1998} show that, for a one-standard-deviation shock $\delta_g =
\sigma_{gg}^{1/2}$ to the innovation of interest, the scaled GIRF is
\begin{equation}
    \psi_Y^g(h)
    \;=\; \sigma_{gg}^{-1/2} \Psi_h \Sigma_u e_g.
    \label{eq:GIRF-Y-closed}
\end{equation}
The contemporaneous response of the remaining variables, $\Sigma_u
e_g/\sigma_{gg}$, depends only on $\Sigma_u$ and not on the ordering of those
variables.\footnote{By Proposition~3.1 of \citet{pesaranShin1998},
\eqref{eq:GIRF-Y-closed} coincides with the orthogonalized impulse response from
a Cholesky decomposition of $\Sigma_u$ in which variable $g$ is ordered first.}
\subsubsection{Propagating GIRFs through lag-selection operators}
\label{sec:lag_section}
The satellite equation uses subsets and lags of the VAR variables. To compact
notation and preserve linearity, we introduce a linear lag-selection operator.
\begin{defn}[\textbf{Linear operator representation with lags}]
\label{def:linear-operator-lags}
\textit{Let $L$ denote the lag operator, with $L^\ell x_t := x_{t-\ell}$
for $\ell \in \mathbb N_0$. For $Y_t \in \R^{n}$ and a maximum lag
$L_{\max} \in \mathbb N_0$, define the polynomial selection operator}
\begin{equation}
    S^{(s)}(L) := \sum_{\ell=0}^{L_{\max}} S^{(s)}_{\ell}\,L^{\ell},
    \qquad S^{(s)}_{\ell} \in \{0,1\}^{m\times n}.
\end{equation}
\textit{The selected vector is the linear transform}
\begin{equation}
    Y_t^{(s)}
    = S^{(s)}(L)\,Y_t
    = \sum_{\ell=0}^{L_{\max}} S^{(s)}_{\ell}\,Y_{t-\ell}
    \in \R^{m}.
\end{equation}
\textit{Each row of the block matrices $\{S^{(s)}_{\ell}\}_{\ell=0}^{L_{\max}}$
contains exactly one entry equal to $1$ across all lags and zeros
elsewhere, thereby selecting exactly one component $(Y_{t-\ell})_j$ per row.
Equivalently, defining the stacked state vector
$\widetilde Y_t = \big( (Y_t)^\top, \ldots, (Y_{t-L_{\max}})^\top \big)^\top$,}
\begin{equation}
    Y_t^{(s)}
    = \underbrace{\big[\,S^{(s)}_{0}\;\;S^{(s)}_{1}\;\;\cdots\;\;S^{(s)}_{L_{\max}}\,\big]}_{=:\,S^{(s)}}
    \,\widetilde Y_t.
\end{equation}
\end{defn}
\medskip
The $g$-th column of each selection matrix $S^{(s)}_\ell$ is zero, so the
variable of interest never enters the satellite regression directly.
\medskip
As an illustration, let
$Y_t = (G_t,\ \mathrm{GDP\;growth}_t,\ \mathrm{unemployment}_t)^\top$
with the variable of interest ordered first, and suppose the satellite relates
$Z_t$ to the current value of GDP growth, its second lag, and the first lag of
unemployment. Setting $n=3$, $m=3$, and $L_{\max}=2$, the block selection
matrices, with the first column set to zero to exclude $G_t$, are
\begin{equation}
    S^{(s)}_{0}=\begin{bmatrix}
    0&1&0\\[2pt] 0&0&0\\[2pt] 0&0&0
    \end{bmatrix},\qquad
    S^{(s)}_{1}=\begin{bmatrix}
    0&0&0\\[2pt] 0&0&0\\[2pt] 0&0&1
    \end{bmatrix},\qquad
    S^{(s)}_{2}=\begin{bmatrix}
    0&0&0\\[2pt] 0&1&0\\[2pt] 0&0&0
    \end{bmatrix},
\end{equation}
and applying the operator yields
\begin{equation}
    Y_t^{(s)}
    = \left( S^{(s)}_{0} + S^{(s)}_{1}L + S^{(s)}_{2}L^2 \right) Y_t
    = \big(\mathrm{GDP\;growth}_{t},\ \mathrm{GDP\;growth}_{t-2},\ \mathrm{unemployment}_{t-1}\big)^\top.
\end{equation}
\medskip
By linearity and time invariance, GIRFs propagate through the selection operator
as a convolution over lags.
\begin{prop}[\textbf{GIRFs for the satellite explanatory vector} $Y^{(s)}$]
    \label{prop:GIRF-Ys}
    Let $Y_t^{(s)} = S^{(s)}(L)\,Y_t$ be the vector of explanatory variables
    entering the satellite equation, where $Y_t\in\R^{n}$ is the VAR vector
    and $S^{(s)}(L)$ is the operator of
    Definition~\ref{def:linear-operator-lags}. For any horizon $h\ge0$,
    history $\omega_{t-1}$, and reduced-form shock $\delta_g\in\R$,
    \begin{equation}
    \label{eq:GIRF-Ys-operator}
        \psi^g_{Y^{(s)}}(h,\delta_g,\omega_{t-1})
        \;=\; S^{(s)}(L)\,\psi^g_{Y}(h,\delta_g,\omega_{t-1})
        \;=\; \sum_{\ell=0}^{L_{\max}} S^{(s)}_\ell\,\psi^g_{Y}(h-\ell,\delta_g,\omega_{t-1}),
    \end{equation}
    with the convention $\psi^g_{Y}(h',\cdot,\cdot)\equiv 0$ for $h'<0$.
\end{prop}
The proof is in Appendix~\ref{app:_Proof_GIRF-Ys}. The GIRF for the systematic
factor follows by applying the same linear operator to the satellite regression.
\begin{prop}[\textbf{GIRF for the systematic factor}]
    \label{prop:GIRF-Z}
    Let $Z_t = \beta_0 + \beta^\top Y_t^{(s)} + \eta_t$ with
    $\E[\eta_t \mid \cdot] = 0$, where $Y_t^{(s)} = S^{(s)}(L)Y_t$. Then,
    for any $h \ge 0$ and $\delta_g \in \R$,
    \begin{equation}
        \psi^g_Z(h,\delta_g,\omega_{t-1})
        =
        \beta^\top S^{(s)}(L)\,
        \psi^g_{Y}(h,\delta_g,\omega_{t-1})
        =
        \sum_{\ell=0}^{L_{\max}}
        \beta^\top S^{(s)}_\ell\,
        \psi^g_{Y}(h-\ell,\delta_g,\omega_{t-1}),
    \end{equation}
    with the convention $\psi^g_{Y}(h',\cdot,\cdot)\equiv 0$ for $h'<0$.
\end{prop}
\begin{proof}
By the satellite equation,
\[
\E[Z_{t+h}\mid u_{gt}=\delta_g,\omega_{t-1}]
= \beta_0 + \beta^\top \E[Y^{(s)}_{t+h}\mid u_{gt}=\delta_g,\omega_{t-1}],
\]
and, similarly,
$\E[Z_{t+h}\mid \omega_{t-1}] = \beta_0 + \beta^\top \E[Y^{(s)}_{t+h}\mid \omega_{t-1}]$,
since the satellite error has zero conditional mean. Subtracting eliminates
the intercept and gives
$\psi^g_Z(h,\delta_g,\omega_{t-1}) = \beta^\top \psi^g_{Y^{(s)}}(h,\delta_g,\omega_{t-1})$.
Applying Proposition~\ref{prop:GIRF-Ys} yields the result.
\end{proof}
\subsubsection{Closed-form generalized impulse responses for portfolio default probabilities}
The final step maps the shock-induced trajectory of the systematic factor into
portfolio default probabilities consistently with the nonlinear Merton structure.
Proposition~\ref{prop:GIRF-Z} delivers the GIRF of $Z$ directly from the GIRF of
the macro-financial variables, hence the full sequence $\{Z_{t+h}\}_{h\ge 0}$
implied by the shock. The transformation $Z \mapsto \pi(Z)$ is nonlinear, so the
PD response is not a linear transformation of $\psi^g_Z(h,\delta_g,\omega_{t-1})$,
and deriving it requires the expectation of a nonlinear transformation of a
Gaussian random variable. We define the GIRF of portfolio default probabilities
as
\begin{equation}
    \psi^g_{\pi(Z)}(h,\delta_g,\omega_{t-1})
    := \E\!\big[\pi(Z_{t+h}) \mid u_{gt}=\delta_g,\ \omega_{t-1}\big]
     - \E\!\big[\pi(Z_{t+h}) \mid \omega_{t-1}\big].
    \label{eq:GIRF_PD_Def}
\end{equation}
The following lemma characterizes the expectation of the Merton mapping under a
Gaussian factor.
\begin{lem}\label{lem:closed}
    Let $W \sim \mathcal N(\mu,\sigma^2)$. Then
    \begin{equation}
        \E\!\left[
            \Phi\!\left(
                \frac{\Phi^{-1}(p)-\sqrt{\rho}\,W}{\sqrt{1-\rho}}
            \right)
        \right]
        =
        \Phi\!\left(
            \frac{\Phi^{-1}(p)-\sqrt{\rho}\,\mu}{\sqrt{(1-\rho)+\rho\,\sigma^2}}
        \right).
        \label{eq:closed}
    \end{equation}
\end{lem}
The proof is in Appendix~\ref{app:proof_lemma}. Let the conditional means and
variances of the unshocked and shocked distributions of $Z_{t+h}$, $h \ge 0$, be
\begin{equation}
    \mu_{t+h} := \E[Z_{t+h} \mid \omega_{t-1}],
    \qquad
    \mu^{(\delta_g)}_{t+h} := \E[Z_{t+h} \mid u_{gt}=\delta_g,\, \omega_{t-1}],
\end{equation}
\begin{equation}
    s_{t+h}^2 := \Var[Z_{t+h} \mid \omega_{t-1}],
    \qquad
    (s^{(\delta_g)}_{t+h})^2 := \Var[Z_{t+h} \mid u_{gt}=\delta_g,\, \omega_{t-1}].
\end{equation}
\begin{prop}[\textbf{Closed-form GIRF of the default probabilities}]\label{prop:GIRF-PD}
    Under the model specification \eqref{eq:var}--\eqref{eq:Z}, the GIRF of
    the point-in-time probability of default $\pi(Z)$ at horizon $h$ in
    response to a shock $\delta_g$ to the innovation of interest is
    \begin{equation}
        \psi^g_{\pi(Z)}(h,\delta_g,\omega_{t-1})
        =
        \Phi\!\left(
            \frac{\Phi^{-1}(p) - \sqrt{\rho}\,\mu_{t+h}^{(\delta_g)}}
            {\sqrt{(1-\rho)+\rho\,(s_{t+h}^{(\delta_g)})^2}}
        \right)
        -
        \Phi\!\left(
            \frac{\Phi^{-1}(p) - \sqrt{\rho}\,\mu_{t+h}}
            {\sqrt{(1-\rho)+\rho\,s_{t+h}^2}}
        \right).
        \label{eq:GIRF_PD}
    \end{equation}
\end{prop}
The proof is in Appendix~\ref{app:proof_Proposition_GIRF_PD}. The expression
depends on the systematic factor through both its conditional mean and its
conditional variance: the shock shifts the mean of $Z_{t+h}$, and, through the
nonlinear Merton map, the conditional dispersion of $Z_{t+h}$ also enters the
level of the default probability. The closed-form GIRF therefore captures more
than a deterministic substitution of the factor path into
\eqref{eq:pd_conditional_z}. Using Definition~\ref{prop:GIRF-GPR}, the shocked
conditional mean is $\mu_{t+h}^{(\delta_g)} = \mu_{t+h} +
\psi^g_{Z}(h,\delta_g,\omega_{t-1})$, so \eqref{eq:GIRF_PD} can be written
directly in terms of the GIRF of the systematic factor:
\begin{equation*}
    \psi^g_{\pi(Z)}(h,\delta_g,\omega_{t-1})
    =
    \Phi\!\left( \frac{\Phi^{-1}(p) - \sqrt{\rho}\,\big(\mu_{t+h} + \psi^g_{Z}(h,\delta_g,\omega_{t-1})\big)}{\sqrt{(1-\rho)+\rho\,(s^{(\delta_g)}_{t+h})^2}} \right)
    -
    \Phi\!\left( \frac{\Phi^{-1}(p) - \sqrt{\rho}\,\mu_{t+h}}{\sqrt{(1-\rho)+\rho\,s_{t+h}^2}} \right).
\end{equation*}
\subsubsection{Conditional mean and variance of the systematic factor}
\label{subsec:mu-s2}
Proposition~\ref{prop:GIRF-PD} shows that the GIRF of portfolio default
probabilities is determined by the conditional mean and variance of the
systematic factor, $(\mu_{t+h}, s_{t+h}^2)$. We compute these quantities using
the lag-selection operator. Throughout, the shock is a scalar innovation to the
equation of the variable of interest, $u_{gt}=\delta_g$.
\begin{prop}[\textbf{Conditional means of the systematic factor}]
\label{prop:Mean}
    Under the system \eqref{eq:var}--\eqref{eq:Z}, the conditional mean of
    $Z_{t+h}$ for any $h\ge 0$ is
    \begin{equation}
    \label{eq:mu-baseline}
        \mu_{t+h}
        =
        \beta_0
        +
        \sum_{\ell=0}^{L_{\max}}
            \beta^\top S^{(s)}_\ell \,\mu^Y_{t+h-\ell},
        \qquad
        \mu^Y_{t+h-\ell}:=\E[Y_{t+h-\ell}\mid\omega_{t-1}].
    \end{equation}
    Under the scalar innovation of interest $u_{gt}=\delta_g$, the shocked
    conditional mean is
    \begin{equation}
        \label{eq:mu-shocked}
        \mu^{(\delta_g)}_{t+h}
        =
        \mu_{t+h}
        +
        \psi_Z^g(h,\delta_g,\omega_{t-1}),
        \qquad
        \psi_Z^g(h,\delta_g,\omega_{t-1})
        =
        \sum_{\ell=0}^{L_{\max}}
            \beta^\top S^{(s)}_\ell\,
            \psi_Y^g(h-\ell,\delta_g,\omega_{t-1}),
    \end{equation}
    with the convention $\psi_Y^g(h',\cdot,\cdot)\equiv 0$ for $h'<0$.
\end{prop}
For the variances, conditioning on $u_{gt}=\delta_g$ fixes only the component of
the contemporaneous innovation associated with the variable of interest; the
remaining components of $u_t$ stay random. Let
\begin{equation}
    \label{eq:sigma-u-cond-g}
    \Sigma_{u\mid g}
    :=
    \Var(u_t\mid u_{gt})
    =
    \Sigma_u
    -
    \frac{\Sigma_u e_g e_g^\top \Sigma_u}{\sigma_{gg}},
    \qquad
    \sigma_{gg}=e_g^\top\Sigma_u e_g.
\end{equation}
\begin{prop}[\textbf{Conditional variances of the systematic factor}]
\label{prop:Var}
    Under the system \eqref{eq:var}--\eqref{eq:Z}, the conditional variance
    of $Z_{t+h}$ for any $h\ge 0$ is
    \begin{equation}
        \label{eq:s2-baseline}
        s^2_{t+h}
        =
        \sum_{q=0}^{h}
            B(h,q)\,\Sigma_u\,B(h,q)^\top
        +
        \sigma_\eta^2.
    \end{equation}
    Under the scalar innovation of interest $u_{gt}=\delta_g$,
    \begin{equation}
        \label{eq:s2-shocked}
        \big(s^{(\delta_g)}_{t+h}\big)^2
        =
        B(h,0)\,\Sigma_{u\mid g}\,B(h,0)^\top
        +
        \sum_{q=1}^{h}
            B(h,q)\,\Sigma_u\,B(h,q)^\top
        +
        \sigma_\eta^2,
    \end{equation}
    where $\sigma_\eta^2$ is the satellite error variance, $\Sigma_u$ the
    VAR innovation covariance, and
    \begin{equation}
        \label{eq:G-hq}
        G_{h,q}
        :=
        \sum_{\ell=0}^{L_{\max}}
            S^{(s)}_\ell\Psi_{h-\ell-q}
        \in \R^{m\times n},
        \qquad
        B(h,q)
        :=
        \beta^\top G_{h,q}
        \in \R^{1\times n},
    \end{equation}
    with the convention $\Psi_r=0$ for $r<0$.
\end{prop}
The proofs of Propositions~\ref{prop:Mean} and~\ref{prop:Var} are in
Appendix~\ref{app:proof_Moments}. Together with Proposition~\ref{prop:GIRF-PD},
they characterize the evolution of portfolio default probabilities in response to
a scalar innovation of interest $\delta_g$ at time $t$. In estimation, these
expressions are evaluated for each joint draw of the VAR parameters
$(c,A_1,\dots,A_p,\Sigma_u)$ and the satellite parameters
$(\beta_0,\beta,\sigma_\eta^2)$, with the satellite draws mixed across
specifications according to the weights $w_s$, and the resulting
default-probability responses are summarized by their posterior medians and
credible bands.
